\chapter{Procedimentos Metodológicos}
\label{cap:metodologia2x}

Apresente aqui uma visão geral do objeto de estudo (base de dados, sistema, população, amostra etc.) e do contexto em que a pesquisa será conduzida. Não podemos deixar entre capítulos e seções, espaço em branco e sem texto, como seria se não existisse esse parágrafo atual, isos vale para o documento todo.

\section{Caracterização da Pesquisa}
\label{sec:caracterizacao}

Classifique a pesquisa quanto às quatro dimensões metodológicas mais comuns, sempre justificando o enquadramento escolhido:

Quanto à natureza, classifique como básica ou aplicada. Quanto à abordagem do problema, classifique como quantitativa, qualitativa ou quali-quantitativa. Quanto aos objetivos, classifique como exploratória, descritiva ou explicativa. Quanto aos procedimentos técnicos, classifique como bibliográfica, documental, experimental, estudo de caso, entre outros \cite{marconi2017metodologia}. Use uma boa referência de livro de metodologia.

\section{Etapas da Pesquisa}

Descreva as macroetapas que estruturam a pesquisa do início ao fim (por exemplo: preparação, definição de variáveis, execução/coleta e análise), explicando brevemente o que ocorre em cada uma antes de detalhá-las em subseções.

\subsection{Etapa 1 -- Nome da Primeira Etapa}

Descreva a primeira etapa da pesquisa (por exemplo: preparação do ambiente, da base de dados ou do material que será utilizado).

\subsection{Etapa 2 -- Definição das Variáveis e do Plano de Testes}

Caso a pesquisa seja experimental, defina aqui a variável de controle (o que será mantido fixo para não interferir na análise), a variável independente (o que será manipulado) e a variável dependente (o que será medido como resultado) \cite{marconi2017metodologia}.

A Tabela \ref{tab:plano_testes} apresenta um exemplo de plano de testes, combinando as variáveis independentes definidas para compor os cenários experimentais.

\begin{table}[H]
\IBGEtab{
\caption{Plano de testes experimental}
\label{tab:plano_testes}
}{
\begin{tabular}{c c c c}
\hline
\textbf{Caso} &
\textbf{Variável A} &
\textbf{Variável B} &
\textbf{Métricas Coletadas} \\
\hline

C1 & Valor 1 & Valor 1 & Métrica 1, Métrica 2 \\
C2 & Valor 1 & Valor 2 & Métrica 1, Métrica 2 \\
C3 & Valor 2 & Valor 1 & Métrica 1, Métrica 2 \\
C4 & Valor 2 & Valor 2 & Métrica 1, Métrica 2 \\

\hline
\end{tabular}
}{
\vspace{0.3cm}
\fonte{Elaborado pelo autor, 2026.}
}
\end{table}

\subsection{Etapa 3 -- Execução e Coleta das Métricas}

Descreva como os experimentos (ou a coleta de dados) serão executados na prática, e com quais ferramentas as métricas de interesse serão obtidas.

\subsection{Etapa 4 -- Análise dos Resultados}

Descreva como os dados coletados serão analisados (por exemplo: estatística descritiva, comparação entre cenários, testes de hipótese) e como essa análise responderá à questão de pesquisa apresentada no Capítulo 2.
